\documentclass[11pt,a4paper]{article}
\usepackage{get-uebung}

\blatt{10}{Ersatzquellen, Leistungsanpassung und fiese Fallen}{Kap.~2.11--2.13}

\begin{document}

\blattkopf{
  \item eine lineare Zweipol-Schaltung durch eine Ersatzspannungs- oder Ersatzstromquelle ($U_0$, $I_k$, $R_i$) ersetzen,
  \item den Innenwiderstand durch Deaktivieren aller Quellen bestimmen,
  \item den Lastwiderstand für maximale Leistung angeben und Leistung sowie Wirkungsgrad bewerten,
  \item Bauelemente ohne Einfluss auf den Rest der Schaltung („fiese Fallen“) sofort erkennen,
  \item den Strom im Brückenzweig einer nicht abgeglichenen Wheatstonebrücke mit einer Ersatzquelle berechnen.
}

% -------------------------------------------------------------------
\begin{aufgabe}{Ersatzspannungs- und Ersatzstromquelle}{16}
Die folgende Schaltung soll bezüglich der Klemmen A--B durch eine Ersatzquelle beschrieben werden.
\begin{center}
\begin{circuitikz}
  \draw (0,0) to[V, v<=$U_q$] (0,3)
        to[R, l=$R_1$] (3,3)
        to[R, l=$R_3$] (6,3) -- (7.5,3) node[ocirc]{} node[right]{A};
  \draw (3,3) to[R, l_=$R_2$, *-*] (3,0);
  \draw (6,0) to[I, i=$I_q$, *-*] (6,3);
  \draw (0,0) -- (7.5,0) node[ocirc]{} node[right]{B};
  \draw (7.5,2.9) to[open, v^=$U_{AB}$] (7.5,0.1);
\end{circuitikz}
\end{center}
\gegeben{$U_q = \SI{24}{V}$, $I_q = \SI{2}{A}$, $R_1 = \SI{4}{\ohm}$, $R_2 = \SI{12}{\ohm}$, $R_3 = \SI{2}{\ohm}$}
\begin{teilaufgaben}
  \item Bestimmen Sie den Innenwiderstand $R_i$ der Schaltung bezüglich der Klemmen A--B.
  \item Berechnen Sie die Leerlaufspannung $U_0 = U_{AB}$ (Klemmen offen).
  \item Geben Sie die Ersatzspannungsquelle und die Ersatzstromquelle mit allen Werten an (Skizze!).
  \item Kontrolle: Berechnen Sie den Kurzschlussstrom $I_k$ direkt in der Originalschaltung, indem Sie A und B kurzschließen.
  \item An die Klemmen wird ein Lastwiderstand $R_L = \SI{9}{\ohm}$ angeschlossen. Wie groß sind Strom und Spannung am Lastwiderstand?
\end{teilaufgaben}
\begin{tipp}Deaktivieren heißt: ideale Spannungsquelle $\to$ Kurzschluss, ideale Stromquelle $\to$ Leerlauf. Für $U_0$ hilft der Überlagerungssatz.\end{tipp}
\end{aufgabe}

\begin{loesung}
\schritt{a) Innenwiderstand}
\leitfrage{Wie sieht die Schaltung aus, wenn $U_q$ durch einen Kurzschluss und $I_q$ durch einen Leerlauf ersetzt wird?}
Der Zweig mit $I_q$ verschwindet. Von den Klemmen aus sieht man $R_3$ in Reihe zur Parallelschaltung von $R_1$ (jetzt nach Masse kurzgeschlossen) und $R_2$:
\[
  R_i = R_3 + R_1\,\|\,R_2 = \SI{2}{\ohm} + \frac{\SI{4}{\ohm}\cdot\SI{12}{\ohm}}{\SI{16}{\ohm}} = \SI{2}{\ohm}+\SI{3}{\ohm} = \ergebnis{\SI{5}{\ohm}}
\]
\schritt{b) Leerlaufspannung mit dem Überlagerungssatz}
\emph{Nur $U_q$ aktiv} ($I_q$ als Leerlauf): Durch $R_3$ fließt kein Strom, also liegt an den Klemmen die Spannung an $R_2$ (Spannungsteiler):
\[
  U_0' = U_q\,\frac{R_2}{R_1+R_2} = \SI{24}{V}\cdot\frac{12}{16} = \SI{18}{V}
\]
\emph{Nur $I_q$ aktiv} ($U_q$ als Kurzschluss): $I_q$ fließt in Klemme A und dann vollständig durch $R_3$ und $R_1\|R_2$ zurück -- also genau durch den Widerstand $R_i$:
\[
  U_0'' = I_q\cdot(R_3 + R_1\|R_2) = \SI{2}{A}\cdot\SI{5}{\ohm} = \SI{10}{V}
\]
\[
  U_0 = U_0' + U_0'' = \ergebnis{\SI{28}{V}}
\]
\schritt{c) Ersatzquellen}
Ersatzspannungsquelle: $U_0 = \SI{28}{V}$ in Reihe mit $R_i = \SI{5}{\ohm}$.\\
Ersatzstromquelle: $I_k = U_0/R_i = \SI{28}{V}/\SI{5}{\ohm} = \ergebnis{\SI{5.6}{A}}$ parallel zu $R_i = \SI{5}{\ohm}$ (Pfeil von $I_k$ zeigt zur Klemme A).
\begin{center}
\begin{circuitikz}
  \draw (0,0) to[V, v<=$U_0$] (0,2.5) to[R, l=$R_i$] (2.5,2.5) to[short,-o] (3,2.5) node[right]{A};
  \draw (0,0) to[short,-o] (3,0) node[right]{B};
  \draw (6,0) to[I, i=$I_k$] (6,2.5) to[short,-o] (8.5,2.5) node[right]{A};
  \draw (7.5,2.5) to[R, l=$R_i$, *-*] (7.5,0);
  \draw (6,0) to[short,-o] (8.5,0) node[right]{B};
\end{circuitikz}
\end{center}
\schritt{d) Kontrolle über den Kurzschlussstrom}
\leitfrage{Bei kurzgeschlossenen Klemmen liegt A auf Masse. Welche Spannung $U_K$ stellt sich am oberen Knoten von $R_2$ ein?}
Knotensatz am Knoten K (Ströme hinein $=$ Ströme heraus), $R_3$ liegt nun zwischen K und Masse:
\[
  \frac{U_q - U_K}{R_1} = \frac{U_K}{R_2} + \frac{U_K}{R_3}
  \;\Rightarrow\; \SI{6}{A} = U_K\left(\tfrac14+\tfrac1{12}+\tfrac12\right)\si{S}
  \;\Rightarrow\; U_K = \SI{7.2}{V}
\]
Im Kurzschluss fließen der Strom durch $R_3$ und der Quellstrom $I_q$ zusammen:
$I_k = U_K/R_3 + I_q = \SI{3.6}{A}+\SI{2}{A} = \SI{5.6}{A}$ \checkmark
\schritt{e) Lastwiderstand anschließen}
Mit der Ersatzspannungsquelle ist das nur noch ein einfacher Stromkreis:
\[
  I_L = \frac{U_0}{R_i+R_L} = \frac{\SI{28}{V}}{\SI{14}{\ohm}} = \ergebnis{\SI{2}{A}},\qquad
  U_L = R_L\cdot I_L = \ergebnis{\SI{18}{V}}
\]
\typfehler{Beim Innenwiderstand die Stromquelle kurzschließen statt sie zu entfernen. Bei $U_0$ den Beitrag der Stromquelle vergessen („durch $R_3$ fließt im Leerlauf kein Strom“ gilt nur, solange $I_q$ deaktiviert ist!).}
\end{loesung}

% -------------------------------------------------------------------
\begin{aufgabe}{Leistungsanpassung}{14}
Eine reale Spannungsquelle ($U_0 = \SI{12}{V}$, $R_i = \SI{2}{\ohm}$) speist einen einstellbaren Lastwiderstand $R_L$.
\begin{center}
\begin{circuitikz}
  \draw (0,0) to[V, v<=$U_0$] (0,2.5) to[R, l=$R_i$] (3,2.5) to[short, i=$I$] (4.5,2.5)
        to[vR, l_=$R_L$, v^=$U$] (4.5,0) -- (0,0);
  \draw[dashed,gray] (-1.1,-0.4) rectangle (3.4,3.25);
\end{circuitikz}
\end{center}
\begin{teilaufgaben}
  \item Bei welchem Wert von $R_L$ nimmt der Lastwiderstand die größte Leistung auf? Wie groß ist $P_\text{max}$?
  \item Vervollständigen Sie die Tabelle. Der Wirkungsgrad ist $\eta = P_L/P_\text{Quelle}$, wobei $P_\text{Quelle} = U_0\cdot I$ die von der idealen Quelle abgegebene Leistung ist.
  \begin{center}
  \renewcommand{\arraystretch}{1.2}
  \begin{tabular}{l|c|c|c|c}
    $R_L$ & $\SI{1}{\ohm}$ & $\SI{2}{\ohm}$ & $\SI{4}{\ohm}$ & $\SI{20}{\ohm}$\\\hline
    $I$ & & & & \\\hline
    $P_L$ & & & & \\\hline
    $\eta$ & & & & \\
  \end{tabular}
  \end{center}
  \item Zeigen Sie allgemein, dass $\eta = \dfrac{R_L}{R_i+R_L}$ gilt. Welcher Wirkungsgrad ergibt sich bei Leistungsanpassung?
  \item Ein Kraftwerk wird nie in Leistungsanpassung betrieben, eine Antenne am Empfänger dagegen schon. Begründen Sie.
  \item Welcher Lastwiderstand entnimmt der Schaltung aus Aufgabe~10.1 die größte Leistung, und wie groß ist diese?
\end{teilaufgaben}
\end{aufgabe}

\begin{loesung}
\schritt{a) Bedingung für maximale Leistung}
Nach Skript Kap.~2.12 ist die Leistung maximal für $R_L = R_i$:
\[
  R_L = \ergebnis{\SI{2}{\ohm}},\qquad P_\text{max} = \frac{U_0^2}{4R_i} = \frac{(\SI{12}{V})^2}{\SI{8}{\ohm}} = \ergebnis{\SI{18}{W}}
\]
\schritt{b) Tabelle}
\leitfrage{Welche drei Formeln brauchen Sie? $I = U_0/(R_i+R_L)$, $P_L = I^2R_L$ und $\eta = P_L/(U_0 I)$.}
\begin{center}
\renewcommand{\arraystretch}{1.2}
\begin{tabular}{l|c|c|c|c}
  $R_L$ & $\SI{1}{\ohm}$ & $\SI{2}{\ohm}$ & $\SI{4}{\ohm}$ & $\SI{20}{\ohm}$\\\hline
  $I$ & $\SI{4}{A}$ & $\SI{3}{A}$ & $\SI{2}{A}$ & $\SI{0.545}{A}$\\\hline
  $P_L$ & $\SI{16}{W}$ & $\SI{18}{W}$ & $\SI{16}{W}$ & $\SI{5.95}{W}$\\\hline
  $\eta$ & $\SI{33.3}{\percent}$ & $\SI{50}{\percent}$ & $\SI{66.7}{\percent}$ & $\SI{90.9}{\percent}$\\
\end{tabular}
\end{center}
Beobachtung: Das Leistungsmaximum liegt bei $R_L = R_i$; der Wirkungsgrad steigt dagegen mit $R_L$ immer weiter an.
\schritt{c) Wirkungsgrad allgemein}
Durch $R_i$ und $R_L$ fließt derselbe Strom $I$:
\[
  \eta = \frac{P_L}{P_\text{Quelle}} = \frac{I^2R_L}{I^2(R_i+R_L)} = \frac{R_L}{R_i+R_L}
  \qquad\xrightarrow{R_L=R_i}\qquad \eta = \ergebnis{\tfrac12 = \SI{50}{\percent}}
\]
Bei Leistungsanpassung wird also genau die Hälfte der Leistung im Innenwiderstand verheizt.
\schritt{d) Kraftwerk vs.\ Antenne}
Im Kraftwerk kommt es auf den \emph{Wirkungsgrad} an: $\SI{50}{\percent}$ Verlust wären unbezahlbar und würden den Generator überhitzen. Deshalb gilt dort $R_L \gg R_i$.
Beim Empfänger ist das Antennensignal winzig, man will \emph{möglichst viel Leistung} herausholen -- der Wirkungsgrad spielt keine Rolle. Deshalb wird dort angepasst.
\schritt{e) Anwendung auf Aufgabe 10.1}
Die Ersatzquelle hat $U_0 = \SI{28}{V}$ und $R_i = \SI{5}{\ohm}$:
\[
  R_L = R_i = \ergebnis{\SI{5}{\ohm}},\qquad P_\text{max} = \frac{(\SI{28}{V})^2}{4\cdot\SI{5}{\ohm}} = \ergebnis{\SI{39.2}{W}}
\]
Genau dafür ist die Ersatzquelle nützlich: Ohne sie wäre diese Frage sehr mühsam.
\typfehler{$P_\text{max} = U_0^2/R_i$ statt $U_0^2/(4R_i)$ -- bei Anpassung liegt nur $U_0/2$ am Lastwiderstand. Leistungsanpassung mit „maximalem Wirkungsgrad“ verwechseln.}
\end{loesung}

% -------------------------------------------------------------------
\begin{aufgabe}{Fiese Fallen}{14}
\textbf{Teil 1:} Gegeben ist die folgende Schaltung.
\begin{center}
\begin{circuitikz}
  \draw (0,0) to[V, v<=$U$, i=$I_U$] (0,3) -- (1.5,3)
        to[R, l=$R_1$, *-*] (4.5,3) -- (7,3)
        to[R, l=$R_5$] (7,1.5) to[I, i<=$I_q$] (7,0) -- (0,0);
  \draw (1.5,3) to[R, l_=$R_4$, -*] (1.5,0);
  \draw (4.5,3) to[R, l=$R_L$, i>^=$I_L$, -*] (4.5,0);
\end{circuitikz}
\end{center}
\gegeben{$U = \SI{10}{V}$, $I_q = \SI{1}{A}$, $R_1 = R_L = \SI{5}{\ohm}$, $R_4 = \SI{100}{\ohm}$, $R_5 = \SI{47}{\ohm}$}
\begin{teilaufgaben}
  \item Welche Widerstände haben keinen Einfluss auf den Strom $I_L$? Begründen Sie kurz.
  \item Berechnen Sie $I_L$.
  \item Berechnen Sie den Strom $I_U$ der Spannungsquelle. Hat $R_4$ wirklich „gar keinen“ Einfluss?
\end{teilaufgaben}
\textbf{Teil 2:} Eine Wheatstonebrücke wird von $U_0 = \SI{9}{V}$ gespeist.
\begin{center}
\begin{circuitikz}
  \draw (0,0) to[V, v<=$U_0$, i=$I_0$] (0,4) -- (6.5,4);
  \draw (2.5,4) to[R, l_=$R_1$, *-*] (2.5,2) to[R, l_=$R_2$, -*] (2.5,0);
  \draw (6.5,4) to[R, l=$R_3$] (6.5,2) to[R, l=$R_4$, *-] (6.5,0) -- (0,0);
  \draw (2.5,2) to[R, l=$R_5$] (6.5,2);
\end{circuitikz}
\end{center}
\gegeben{$R_1 = \SI{100}{\ohm}$, $R_2 = \SI{200}{\ohm}$, $R_3 = \SI{150}{\ohm}$, $R_4 = \SI{300}{\ohm}$, $R_5 = \SI{1}{k\ohm}$}
\begin{teilaufgaben}[resume]
  \item Berechnen Sie den Strom $I_0$, ohne ein Gleichungssystem aufzustellen.
  \item Wie ändert sich $I_0$, wenn $R_5$ durch einen Kurzschluss ersetzt wird?
\end{teilaufgaben}
\end{aufgabe}

\begin{loesung}
\schritt{a) Einflussfreie Bauelemente erkennen}
\leitfrage{Welcher Widerstand liegt direkt parallel zu einer idealen Spannungsquelle, welcher in Reihe zu einer idealen Stromquelle?}
\begin{itemize}[nosep]
  \item $R_4$ liegt parallel zur idealen Spannungsquelle: An ihm liegt immer $U$, egal wie groß er ist. Für den Rest der Schaltung bleibt die Spannung am linken Knoten $U$.
  \item $R_5$ liegt in Reihe zur idealen Stromquelle: Durch ihn fließt immer $I_q$, der Zweig liefert unabhängig von $R_5$ den Strom $I_q$ in den Knoten.
\end{itemize}
Also: $\ergebnis{R_4 \text{ und } R_5}$ haben keinen Einfluss auf $I_L$. Wir dürfen $R_4$ weglassen und $R_5$ kurzschließen.
\schritt{b) Laststrom}
Knotensatz am oberen Knoten von $R_L$ mit dem Potential $U_A$:
\[
  \frac{U-U_A}{R_1} + I_q = \frac{U_A}{R_L}
  \;\Rightarrow\; \SI{2}{A} - \frac{U_A}{\SI{5}{\ohm}} + \SI{1}{A} = \frac{U_A}{\SI{5}{\ohm}}
  \;\Rightarrow\; U_A = \SI{7.5}{V}
\]
\[
  I_L = \frac{U_A}{R_L} = \ergebnis{\SI{1.5}{A}}
\]
\schritt{c) Strom der Spannungsquelle}
Die Quelle versorgt $R_4$ und $R_1$:
\[
  I_U = \frac{U}{R_4} + \frac{U-U_A}{R_1} = \SI{0.1}{A} + \SI{0.5}{A} = \ergebnis{\SI{0.6}{A}}
\]
$R_4$ beeinflusst also sehr wohl den Strom (und die Leistung) der \emph{Quelle selbst} -- nur der \emph{Rest} der Schaltung merkt nichts davon. Genauso hängt die Spannung an der Stromquelle von $R_5$ ab.
\schritt{d) Abgeglichene Brücke}
\leitfrage{Ist die Brücke abgeglichen? Prüfen Sie $R_1/R_2 = R_3/R_4$.}
$R_1/R_2 = 100/200 = 0{,}5$ und $R_3/R_4 = 150/300 = 0{,}5$ \checkmark. Die Brücke ist abgeglichen, durch $R_5$ fließt kein Strom -- wir dürfen ihn durch einen Leerlauf ersetzen:
\[
  R_\text{ges} = (R_1+R_2)\,\|\,(R_3+R_4) = \frac{\SI{300}{\ohm}\cdot\SI{450}{\ohm}}{\SI{750}{\ohm}} = \SI{180}{\ohm},
  \qquad I_0 = \frac{\SI{9}{V}}{\SI{180}{\ohm}} = \ergebnis{\SI{50}{mA}}
\]
\schritt{e) $R_5$ kurzgeschlossen}
Gar nicht: $I_0 = \SI{50}{mA}$. Kontrolle: Mit Kurzschluss ist $R_\text{ges} = R_1\|R_3 + R_2\|R_4 = \SI{60}{\ohm} + \SI{120}{\ohm} = \SI{180}{\ohm}$ \checkmark
\typfehler{Mit großem Aufwand ein Gleichungssystem aufstellen, statt zuerst nach einflussfreien Bauelementen zu suchen. Aus „kein Einfluss auf den Rest“ fälschlich schließen, dass auch der Quellenstrom unverändert bleibt.}
\end{loesung}

% -------------------------------------------------------------------
\begin{aufgabe}{Widerstand im Brückenzweig}{14}
In der nicht abgeglichenen Brücke soll der Strom $I_5$ im Brückenwiderstand $R_5$ bestimmt werden.
\begin{center}
\begin{circuitikz}
  \draw (0,0) to[V, v<=$U_0$] (0,4) -- (6.5,4);
  \draw (2.5,4) to[R, l_=$R_1$, *-*] (2.5,2) node[left=2pt]{A} to[R, l_=$R_2$, -*] (2.5,0);
  \draw (6.5,4) to[R, l=$R_3$] (6.5,2) node[right=2pt]{B} to[R, l=$R_4$, *-] (6.5,0) -- (0,0);
  \draw (2.5,2) to[R, l=$R_5$, i>_=$I_5$] (6.5,2);
\end{circuitikz}
\end{center}
\gegeben{$U_0 = \SI{12}{V}$, $R_1 = R_2 = \SI{2}{k\ohm}$, $R_3 = \SI{1}{k\ohm}$, $R_4 = \SI{3}{k\ohm}$, $R_5 = \SI{1.25}{k\ohm}$}
\begin{teilaufgaben}
  \item Trennen Sie $R_5$ heraus und berechnen Sie die Leerlaufspannung $U_{AB,0}$ zwischen den Punkten A und B.
  \item Bestimmen Sie den Innenwiderstand $R_i$ der Restschaltung zwischen A und B.
  \item Berechnen Sie $I_5$ und die Spannung an $R_5$. In welche Richtung fließt der Strom tatsächlich?
  \item Welcher Wert von $R_5$ würde die größte Leistung aufnehmen? Wie groß wäre sie?
\end{teilaufgaben}
\end{aufgabe}

\begin{loesung}
\schritt{a) Leerlaufspannung}
\leitfrage{Ohne $R_5$ bleiben zwei unabhängige Spannungsteiler. Welche Spannung liegt jeweils am unteren Widerstand?}
\[
  U_{A0} = U_0\frac{R_2}{R_1+R_2} = \SI{12}{V}\cdot\frac{2}{4} = \SI{6}{V},\qquad
  U_{B0} = U_0\frac{R_4}{R_3+R_4} = \SI{12}{V}\cdot\frac{3}{4} = \SI{9}{V}
\]
Maschensatz über die beiden unteren Widerstände:
\[
  U_{AB,0} = U_{A0} - U_{B0} = \ergebnis{-\SI{3}{V}}
\]
\schritt{b) Innenwiderstand}
Die Spannungsquelle wird zum Kurzschluss -- oberer und unterer Draht liegen dann auf demselben Potential. Von A aus liegen $R_1$ und $R_2$ parallel, von B aus $R_3$ und $R_4$; beide Parallelschaltungen liegen in Reihe:
\[
  R_i = R_1\|R_2 + R_3\|R_4 = \SI{1}{k\ohm} + \SI{0.75}{k\ohm} = \ergebnis{\SI{1.75}{k\ohm}}
\]
\schritt{c) Strom im Brückenzweig}
Ersatzspannungsquelle ($U_{AB,0}$, $R_i$) mit $R_5$ als Last:
\[
  I_5 = \frac{U_{AB,0}}{R_i+R_5} = \frac{-\SI{3}{V}}{\SI{3}{k\ohm}} = \ergebnis{-\SI{1}{mA}},\qquad
  U_5 = R_5 I_5 = \ergebnis{-\SI{1.25}{V}}
\]
Das negative Vorzeichen bedeutet: Der Strom fließt tatsächlich von B nach A -- plausibel, denn B liegt im Leerlauf auf dem höheren Potential.
\schritt{d) Maximale Leistung}
Leistungsanpassung: $R_5 = R_i = \ergebnis{\SI{1.75}{k\ohm}}$,
\[
  P_\text{max} = \frac{U_{AB,0}^2}{4R_i} = \frac{(\SI{3}{V})^2}{\SI{7}{k\ohm}} \approx \ergebnis{\SI{1.29}{mW}}
\]
\typfehler{Beim Innenwiderstand $R_1+R_2$ parallel zu $R_3+R_4$ rechnen -- das wäre der Widerstand, den die \emph{Quelle} sieht, nicht der zwischen A und B. Das Vorzeichen von $U_{AB,0}$ unterschlagen und so die Stromrichtung falsch angeben.}
\end{loesung}

\end{document}
